% Kuis 1 --- Logika Proposisional
% Satu soal PG per bahasan di ppt/Logika_Preposisi.tex (bukan dari Latihan Soal).
\documentclass[12pt,a4paper]{article}

\usepackage[utf8]{inputenc}
\usepackage[T1]{fontenc}
\usepackage[indonesian]{babel}

\usepackage{geometry}
\geometry{margin=2cm}

\usepackage{lmodern}
\usepackage{amsmath,amssymb}
\usepackage{enumitem}
\usepackage{tabularx}
\usepackage{microtype}

\usepackage{fancyhdr}
\usepackage{lastpage}

\setlength{\headheight}{14pt}
\addtolength{\topmargin}{-2pt}
\usepackage[hidelinks,breaklinks]{hyperref}

\newcommand{\JudulKuis}{Kuis 1}
\newcommand{\KodeMK}{TIF-xxx}
\newcommand{\NamaMK}{Logika Matematika}
\newcommand{\MateriKuis}{Logika Proposisional}
\newcommand{\WaktuKuis}{90 menit}
\newcommand{\NilaiMaksimal}{100}
\newcommand{\Prodi}{Teknik Informatika}
\newcommand{\Fakultas}{Fakultas Teknologi Informasi}
\newcommand{\Institusi}{Universitas Bale Bandung}
\newcommand{\DosenPengampu}{[Nama Dosen Pengampu]}
\newcommand{\TanggalKuis}{[DD Bulan YYYY]}

\pagestyle{fancy}
\fancyhf{}
\fancyhead[L]{\small \JudulKuis\ --- \NamaMK}
\fancyhead[R]{\small \MateriKuis}
\fancyfoot[C]{\small Halaman \thepage\ dari \pageref*{LastPage}}
\renewcommand{\headrulewidth}{0.4pt}
\renewcommand{\footrulewidth}{0pt}

\setlength{\parindent}{0pt}
\setlength{\parskip}{0.5em}

\newlist{pilihan}{enumerate}{1}
\setlist[pilihan]{label=\textbf{\Alph*.}, leftmargin=2.2em, align=left, itemsep=0.35em, topsep=0.35em}

\newcommand{\kuncijawaban}[2]{%
  \par\vspace{0.25em}%
  {\small
  \textit{\textbf{Kunci jawaban:} #1}\par
  \textit{\textbf{Penjelasan:} #2}%
  \par\vspace{0.55em}}%
}

\newcounter{soalkuis}
\newcommand{\soalkuis}{%
  \refstepcounter{soalkuis}%
  \par\vspace{0.6ex plus 0.15ex}%
  \noindent\textbf{Soal \thesoalkuis.}\quad
}

\begin{document}

\begin{center}
  {\Large\bfseries \JudulKuis}\\[0.3em]
  {\large \NamaMK\ (\KodeMK)}\\[0.25em]
  {\normalsize Materi: \MateriKuis}\\[0.35em]
  {\normalsize\textit{(Lembar kunci jawaban dan pembahasan)}}\\[0.8em]
\end{center}

\noindent
\begin{tabularx}{\linewidth}{@{}p{4.5cm} @{\hspace{0.35em}:\hspace{0.65em}} >{\raggedright\arraybackslash}X@{}}
  Nama Mahasiswa        & \rule{0.72\linewidth}{0.4pt} \\
  Nomor Induk Mahasiswa & \rule{0.72\linewidth}{0.4pt} \\
  Kelas / Kelompok      & \rule{0.72\linewidth}{0.4pt} \\
  Program Studi         & \Prodi \\
  Fakultas              & \Fakultas \\
  Institusi             & \Institusi \\
  Dosen Pengampu        & \DosenPengampu \\
  Tanggal kuis          & \TanggalKuis \\
  Waktu pengerjaan      & \WaktuKuis \\
  Nilai maksimal        & \NilaiMaksimal \\
\end{tabularx}

\vspace{0.8em}
\hrule
\vspace{0.8em}

\section*{Petunjuk}
\begin{itemize}[leftmargin=*, itemsep=0.25em]
  \item Kuis ini terdiri atas \textbf{22 soal pilihan ganda} (satu jawaban benar per soal), mencakup seluruh bahasan \textit{Logika Proposisional} pada materi perkuliahan: pengertian logika, kalimat tertutup/terbuka, nilai kebenaran, lambang pernyataan dan \(\tau\), operator \(\neg\), \(\wedge\), \(\vee\), \(\rightarrow\), \(\leftrightarrow\), varian implikasi, tautologi, negasi pernyataan majemuk, hirarki operator, ekuivalen logis, hukum logika, serta penyederhanaan ekspresi.
  \item Isi \textbf{nama, NIM, dan kelas} pada bagian identitas di atas sebelum mengerjakan.
  \item Bacalah teks soal dan label bahasan (cetak miring di awal soal) untuk memastikan konteks pertanyaan.
  \item Pilih \textbf{satu} jawaban terbaik; lingkari atau tulis huruf pilihan (A, B, C, atau D) dengan jelas pada lembar jawaban.
  \item Soal bertipe konsep dan perhitungan singkat nilai kebenaran: gunakan pengertian tabel kebenaran, hukum De Morgan, negasi implikasi (\(p \wedge \neg q\)), kontraposisi, serta hirarki operator sesuai slide kuliah---tanpa perlu menuliskan tabel lengkap kecuali diminta secara eksplisit.
  \item Perangkat elektronik, catatan, dan komunikasi antar peserta \textbf{tidak diperbolehkan} kecuali diizinkan pengawas.
  \item Lembar kuis \textbf{wajib dikumpulkan} kepada pengawas segera setelah waktu \WaktuKuis\ habis; keterlambatan mengikuti aturan kelas.
\end{itemize}

\section*{Soal}

% --- Definisi Logika ---
\soalkuis
\textit{(Definisi Logika)}\\
Menurut pengertian logika matematika dalam materi kuliah, bidang ini terutama digunakan untuk \ldots
\begin{pilihan}
  \item Menghitung luas bangun datar tanpa bukti
  \item Menentukan nilai kebenaran pernyataan dan penarikan kesimpulan menurut aturan yang berlaku
  \item Menerjemahkan teks sastra ke puisi
  \item Mengganti semua argumen dengan opini pribadi
\end{pilihan}
\kuncijawaban{B}{Logika matematika mempelajari penentuan nilai kebenaran pernyataan serta penarikan kesimpulan berdasarkan aturan logika yang berlaku, bukan sekadar perhitungan geometri atau opini subjektif.}

% --- Kalimat Tertutup ---
\soalkuis
\textit{(Kalimat Tertutup / Preposisi)}\\
Kalimat tertutup (preposisi) didefinisikan sebagai kalimat yang \ldots
\begin{pilihan}
  \item Bernilai benar dan salah sekaligus pada situasi yang sama
  \item Hanya berupa pertanyaan retoris
  \item Bernilai benar saja atau salah saja (tidak keduanya sekaligus)
  \item Selalu mengandung kata tanya
\end{pilihan}
\kuncijawaban{C}{Preposisi (kalimat tertutup) hanya boleh bernilai benar atau salah, tidak boleh keduanya pada konteks yang sama; ini ciri kalimat deklaratif yang dapat ditentukan kebenarannya.}

% --- Kalimat Terbuka ---
\soalkuis
\textit{(Kalimat Terbuka)}\\
Manakah yang merupakan contoh \textbf{kalimat terbuka}?
\begin{pilihan}
  \item \(7\) adalah bilangan prima
  \item \(n + 4 > 10\)
  \item Semua mamalia bernapas
  \item \(2 \times 3 = 6\)
\end{pilihan}
\kuncijawaban{B}{Kalimat \(n + 4 > 10\) memuat variabel \(n\); nilai kebenarannya belum tetap sehingga termasuk kalimat terbuka, berbeda dari pernyataan tanpa peubah.}

% --- Nilai Kebenaran ---
\soalkuis
\textit{(Nilai Kebenaran Suatu Pernyataan)}\\
Menentukan kebenaran pernyataan ``Ibukota Australia adalah Sydney'' dengan memeriksa fakta geografis terkini termasuk dalam pendekatan \ldots
\begin{pilihan}
  \item Dasar empiris
  \item Dasar tak empiris semata (hanya simbol)
  \item Hukum De Morgan
  \item Kontraposisi
\end{pilihan}
\kuncijawaban{A}{Pernyataan tentang ibukota diuji dengan fakta di dunia nyata; itu contoh penentuan kebenaran secara \textbf{dasar empiris}, bukan hanya manipulasi simbol.}

% --- Lambang Pernyataan ---
\soalkuis
\textit{(Lambang Suatu Pernyataan)}\\
Penulisan \(r\) : ``Cuaca hari ini cerah'' mengikuti konvensi bahwa huruf kecil \(r\) berfungsi sebagai \ldots
\begin{pilihan}
  \item Lambang untuk seluruh kalimat pernyataan tersebut
  \item Lambang nilai kebenaran \(\tau\)
  \item Operator biimplikasi
  \item Variabel bebas yang tidak boleh diganti
\end{pilihan}
\kuncijawaban{A}{Notasi \(r\) : ``\ldots'' memakai huruf kecil sebagai \textbf{lambang pernyataan} (bukan \(\tau\) yang khusus untuk nilai kebenaran).}

% --- Lambang tau ---
\soalkuis
\textit{(Lambang Nilai Kebenaran)}\\
Jika diketahui \(\tau(p) = F\), maka pernyataan yang dilambangkan \(p\) adalah \ldots
\begin{pilihan}
  \item Pernyataan benar
  \item Pernyataan salah
  \item Kalimat terbuka
  \item Tautologi
\end{pilihan}
\kuncijawaban{B}{\(\tau(p) = F\) berarti pernyataan \(p\) \textbf{salah} (nilai kebenaran False), sesuai pembacaan lambang \(\tau\) pada materi.}

% --- Pernyataan Majemuk ---
\soalkuis
\textit{(Pernyataan Majemuk)}\\
Dalam tabel kata sambung logika, simbol \(\leftrightarrow\) sesuai dengan kata sambung \ldots
\begin{pilihan}
  \item Dan
  \item Atau
  \item Jika dan hanya jika
  \item Tidak
\end{pilihan}
\kuncijawaban{C}{Simbol \(\leftrightarrow\) adalah \textbf{biimplikasi}, yaitu kata sambung ``jika dan hanya jika'' dalam tabel operator logika.}

% --- Konjungsi ---
\soalkuis
\textit{(Nilai Kebenaran Konjungsi)}\\
Jika \(\tau(p) = T\) dan \(\tau(q) = F\), maka \(\tau(p \wedge q)\) adalah \ldots
\begin{pilihan}
  \item \(T\)
  \item \(F\)
  \item Tidak dapat ditentukan
  \item Selalu \(T\)
\end{pilihan}
\kuncijawaban{B}{Konjungsi \(p \wedge q\) benar hanya jika kedua komponen benar; dengan \(\tau(p)=T\) dan \(\tau(q)=F\) diperoleh \(\tau(p \wedge q)=F\).}

% --- Disjungsi ---
\soalkuis
\textit{(Nilai Kebenaran Disjungsi)}\\
Jika \(\tau(p) = F\) dan \(\tau(q) = F\), maka \(\tau(p \vee q)\) adalah \ldots
\begin{pilihan}
  \item \(T\)
  \item \(F\)
  \item \(T\) hanya jika \(p\) dan \(q\) ekuivalen
  \item Sama dengan \(\tau(p \rightarrow q)\)
\end{pilihan}
\kuncijawaban{B}{Disjungsi \(p \vee q\) salah hanya jika \(p\) dan \(q\) keduanya salah; jadi \(\tau(p \vee q)=F\) ketika keduanya \(F\).}

% --- Negasi ---
\soalkuis
\textit{(Nilai Kebenaran Negasi)}\\
Jika \(\tau(p) = F\), maka \(\tau(\neg p)\) adalah \ldots
\begin{pilihan}
  \item \(F\)
  \item \(T\)
  \item \(F\) atau \(T\) bergantung konteks
  \item Sama dengan \(\tau(p)\)
\end{pilihan}
\kuncijawaban{B}{Negasi membalik nilai kebenaran: jika \(\tau(p)=F\) maka \(\tau(\neg p)=T\).}

% --- Implikasi ---
\soalkuis
\textit{(Nilai Kebenaran Implikasi)}\\
Implikasi \(p \rightarrow q\) bernilai \textbf{salah} pada baris tabel kebenaran ketika \ldots
\begin{pilihan}
  \item \(p\) salah dan \(q\) benar
  \item \(p\) benar dan \(q\) salah
  \item \(p\) salah dan \(q\) salah
  \item \(p\) benar dan \(q\) benar
\end{pilihan}
\kuncijawaban{B}{Pada tabel implikasi, satu-satunya baris yang menghasilkan nilai salah adalah anteceden benar (\(p\) benar) dan konsekuen salah (\(q\) salah).}

% --- Biimplikasi ---
\soalkuis
\textit{(Nilai Kebenaran Biimplikasi)}\\
Jika \(\tau(p) = T\) dan \(\tau(q) = F\), maka \(\tau(p \leftrightarrow q)\) adalah \ldots
\begin{pilihan}
  \item \(T\)
  \item \(F\)
  \item \(T\) karena \(p\) benar
  \item Sama dengan \(\tau(p \rightarrow q)\)
\end{pilihan}
\kuncijawaban{B}{Biimplikasi benar hanya jika \(p\) dan \(q\) sama-sama benar atau sama-sama salah; kombinasi \(T,F\) memberikan \(\tau(p \leftrightarrow q)=F\).}

% --- Varian Implikasi ---
\soalkuis
\textit{(Varian dari Implikasi)}\\
Kontraposisi dari implikasi \(p \rightarrow q\) adalah \ldots
\begin{pilihan}
  \item \(q \rightarrow p\)
  \item \(\neg p \rightarrow \neg q\)
  \item \(\neg q \rightarrow \neg p\)
  \item \(p \leftrightarrow q\)
\end{pilihan}
\kuncijawaban{C}{Kontraposisi dari \(p \rightarrow q\) adalah \(\neg q \rightarrow \neg p\); bentuk ini ekuivalen dengan implikasi asli, berbeda dari konvers atau invers.}

% --- Tautologi dan Kontradiksi ---
\soalkuis
\textit{(Tautologi dan Kontradiksi)}\\
Pernyataan majemuk \(p \vee \neg p\) termasuk jenis \ldots
\begin{pilihan}
  \item Tautologi
  \item Kontradiksi
  \item Kontingensi saja (bukan tautologi)
  \item Kalimat terbuka
\end{pilihan}
\kuncijawaban{A}{\(p \vee \neg p\) selalu benar untuk semua nilai \(p\), sehingga termasuk \textbf{tautologi} (hukum eksklusi ketiga).}

% --- Negasi Konjungsi ---
\soalkuis
\textit{(Negasi Konjungsi)}\\
Bentuk yang ekuivalen dengan \(\neg(p \wedge q)\) menurut hukum De Morgan adalah \ldots
\begin{pilihan}
  \item \(\neg p \wedge \neg q\)
  \item \(\neg p \vee \neg q\)
  \item \(p \vee q\)
  \item \(\neg p \rightarrow \neg q\)
\end{pilihan}
\kuncijawaban{B}{Hukum De Morgan untuk konjungsi: \(\neg(p \wedge q) \equiv \neg p \vee \neg q\).}

% --- Negasi Disjungsi ---
\soalkuis
\textit{(Negasi Disjungsi)}\\
Negasi dari \(p \vee q\) dapat ditulis sebagai \ldots
\begin{pilihan}
  \item \(\neg p \vee \neg q\)
  \item \(\neg p \wedge \neg q\)
  \item \(p \wedge q\)
  \item \(\neg(p \wedge q)\)
\end{pilihan}
\kuncijawaban{B}{Hukum De Morgan untuk disjungsi: \(\neg(p \vee q) \equiv \neg p \wedge \neg q\).}

% --- Negasi Implikasi ---
\soalkuis
\textit{(Negasi Implikasi)}\\
Negasi dari \(p \rightarrow q\) secara logis sama dengan \ldots
\begin{pilihan}
  \item \(\neg p \rightarrow \neg q\)
  \item \(p \wedge \neg q\)
  \item \(\neg p \vee \neg q\)
  \item \(q \rightarrow p\)
\end{pilihan}
\kuncijawaban{B}{Negasi implikasi tidak sama dengan invers; \(\neg(p \rightarrow q) \equiv p \wedge \neg q\) (anteceden benar tetapi konsekuen salah).}

% --- Negasi Biimplikasi ---
\soalkuis
\textit{(Negasi Biimplikasi)}\\
Salah satu bentuk yang ekuivalen dengan \(\neg(p \leftrightarrow q)\) adalah \ldots
\begin{pilihan}
  \item \((p \wedge \neg q) \vee (q \wedge \neg p)\)
  \item \(p \leftrightarrow \neg q\) saja, tanpa disjungsi
  \item \(\neg p \leftrightarrow \neg q\)
  \item \(p \rightarrow q\)
\end{pilihan}
\kuncijawaban{A}{\(\neg(p \leftrightarrow q)\) ekuivalen dengan ``salah satu benar, yang lain salah'', yaitu \((p \wedge \neg q) \vee (q \wedge \neg p)\), sesuai materi negasi biimplikasi.}

% --- Hirarki Operator ---
\soalkuis
\textit{(Hirarki Operator Logika)}\\
Menurut hirarki pada materi kuliah, operator yang \textbf{pertama kali} dievaluasi (prioritas tertinggi) adalah \ldots
\begin{pilihan}
  \item \(\neg\) (negasi)
  \item \(\wedge\) (konjungsi)
  \item \(\vee\) (disjungsi)
  \item \(\rightarrow\) (implikasi)
\end{pilihan}
\kuncijawaban{A}{Hirarki operator: \textbf{negasi} (\(\neg\)) dievaluasi lebih dulu (prioritas tertinggi), baru konjungsi, disjungsi, implikasi, dan biimplikasi.}

% --- Ekuivalen Logis ---
\soalkuis
\textit{(Ekuivalen Logis)}\\
Dua preposisi \(p\) dan \(q\) disebut ekuivalen logis (\(p \equiv q\)) jika \ldots
\begin{pilihan}
  \item Keduanya memakai huruf proposisional yang sama
  \item \(\tau(p) = \tau(q)\) pada setiap kombinasi nilai kebenaran
  \item Keduanya merupakan tautologi
  \item \(p \rightarrow q\) selalu benar
\end{pilihan}
\kuncijawaban{B}{Dua preposisi ekuivalen jika kolom nilai kebenaran mereka sama pada \textbf{setiap} baris tabel kebenaran (\(\tau(p)=\tau(q)\) untuk semua interpretasi).}

% --- Hukum Logika ---
\soalkuis
\textit{(Hukum-hukum Logika)}\\
Menurut hukum involusi (negasi ganda), ekspresi \(\neg(\neg p)\) ekuivalen dengan \ldots
\begin{pilihan}
  \item \(\neg p\)
  \item \(p\)
  \item \(T\)
  \item \(F\)
\end{pilihan}
\kuncijawaban{B}{Hukum involusi menyatakan bahwa negasi ganda mengembalikan pernyataan semula: \(\neg(\neg p) \equiv p\).}

% --- Penyederhanaan ---
\soalkuis
\textit{(Penyederhanaan Ekspresi Logika)}\\
Penyederhanaan ekspresi logika pada materi kuliah dilakukan dengan menerapkan \ldots
\begin{pilihan}
  \item Hukum-hukum logika secara berurutan hingga bentuk lebih sederhana
  \item Hanya tabel kebenaran tanpa aturan
  \item Mengabaikan tanda kurung
  \item Mengganti semua \(\vee\) menjadi \(\wedge\) secara acak
\end{pilihan}
\kuncijawaban{A}{Penyederhanaan dilakukan dengan menerapkan hukum-hukum logika (identitas, De Morgan, distributif, absorpsi, dan lain-lain) secara sistematis hingga bentuk ekspresi lebih ringkas.}

\end{document}
